%! TEX program = xelatex
\documentclass[10pt]{book}
\usepackage[utf8]{inputenc}
\usepackage{fontspec}
\usepackage[margin=0.5in]{geometry}
\usepackage{xcolor}
\usepackage{listings}
\usepackage{titlesec}
\usepackage{fancyhdr}
\usepackage{tocloft}

% Set fonts for e-reader
\setmainfont{QTBookmann}
\setmonofont{VictorMono Nerd Font}[
    SizeFeatures={Size=8},
    Color=000000
]

% High contrast colors for e-reader
\definecolor{codebg}{RGB}{255,255,255}
\definecolor{codeframe}{RGB}{0,0,0}
\definecolor{keyword}{RGB}{0,0,139}
\definecolor{comment}{RGB}{139,0,0}
\definecolor{string}{RGB}{0,100,0}
\definecolor{number}{RGB}{0,0,0}

% Configure listings
\lstdefinestyle{ereader}{
    backgroundcolor=\color{codebg},
    frame=single,
    framesep=3pt,
    framerule=1pt,
    rulecolor=\color{codeframe},
    basicstyle=\ttfamily\footnotesize,
    keywordstyle=\color{keyword}\bfseries,
    commentstyle=\color{comment},
    stringstyle=\color{string},
    numberstyle=\color{number},
    breakatwhitespace=false,
    breaklines=true,
    captionpos=b,
    keepspaces=true,
    numbers=left,
    numbersep=5pt,
    showspaces=false,
    showstringspaces=false,
    showtabs=false,
    tabsize=2
}

\lstset{style=ereader}

% Define JSON language
\lstdefinelanguage{JSON}{
    basicstyle=\ttfamily\footnotesize,
    numbers=left,
    numberstyle=\color{number},
    frame=single,
    framesep=3pt,
    framerule=1pt,
    rulecolor=\color{codeframe},
    stringstyle=\color{string},
    showstringspaces=false,
    breaklines=true,
    literate=
     *{0}{{{\color{number}0}}}{1}
      {1}{{{\color{number}1}}}{1}
      {2}{{{\color{number}2}}}{1}
      {3}{{{\color{number}3}}}{1}
      {4}{{{\color{number}4}}}{1}
      {5}{{{\color{number}5}}}{1}
      {6}{{{\color{number}6}}}{1}
      {7}{{{\color{number}7}}}{1}
      {8}{{{\color{number}8}}}{1}
      {9}{{{\color{number}9}}}{1}
      {:}{{{\color{keyword}{:}}}}{1}
      {,}{{{\color{keyword}{,}}}}{1}
      {\{}{{{\color{keyword}{\{}}}}{1}
      {\}}{{{\color{keyword}{\}}}}}{1}
      {[}{{{\color{keyword}{[}}}}{1}
      {]}{{{\color{keyword}{]}}}}{1},
}

% Chapter formatting
\titleformat{\chapter}[display]
{\normalfont\LARGE\bfseries}{\chaptertitlename\ \thechapter}{0pt}{\Huge}
\titleformat{\section}
{\normalfont\Large\bfseries}{\thesection}{1em}{}
\titleformat{\subsection}
{\normalfont\large\bfseries}{\thesubsection}{1em}{}

% Page layout
\pagestyle{fancy}
\fancyhf{}
\fancyhead[LE,RO]{\thepage}
\fancyhead[RE]{\leftmark}
\fancyhead[LO]{\rightmark}

\begin{document}

\title{\Huge\textbf{LovdogModelSelector}\\ \Large Automated ML Framework}
\author{}
\date{}
\maketitle

\tableofcontents

\chapter{Data Management Core (LovdogDF.py)}

The \textbf{LovdogDataFrames} class is the central data management system that orchestrates all data loading, organization, and preprocessing. It serves as the foundation upon which feature extraction and model selection are built.

\section{Main Class: LovdogDataFrames}

\begin{lstlisting}[language=Python]
class LovdogDataFrames:
    def __init__(self, json_file_path=None, 
                 loader_function=None):
        self.json_file_path = json_file_path
        self.loader_function = loader_function
        self.json_data = None
        self.frames = {}
        self.class_names = []
        
        if json_file_path:
            self.load_json(json_file_path)
\end{lstlisting}

\textbf{Explanation:}
The constructor initializes the data management system with two key parameters:
\begin{itemize}
\item \texttt{json\_file\_path}: Path to a JSON configuration file that defines your data sources, classes, and loading parameters
\item \texttt{loader\_function}: A custom function that handles the actual data loading from files (defaults to a simple path returner)
\end{itemize}

The class maintains several important data structures:
\begin{itemize}
\item \texttt{frames}: A dictionary where keys are class names and values are lists of loaded data (images, arrays, etc.)
\item \texttt{class\_names}: List of all unique class labels found in the data
\item \texttt{extra\_data}: Stores additional metadata for each class
\end{itemize}

\section{JSON Configuration System}

\begin{lstlisting}[language=JSON]
{
    "csv": "Features_Table.csv",
    "class1": {
        "class": "positive",
        "path": ["data/class1/"],
        "extra": {
            "cols": 100, 
            "rows": 100, 
            "type": "2d"
        }
    },
    "class2": {
        "class": "negative", 
        "path": ["data/class2/"],
        "extra": {
            "cols": 100, 
            "rows": 100, 
            "type": "2d"
        }
    }
}
\end{lstlisting}

\textbf{Explanation:}
The JSON configuration provides a declarative way to define your entire dataset:
\begin{itemize}
\item \texttt{csv}: Optional name for the output features table
\item Each class definition contains:
  \begin{itemize}
  \item \texttt{class}: The label for this category
  \item \texttt{path}: List of directories or files containing data
  \item \texttt{extra}: Loading parameters specific to your data format
  \end{itemize}
\end{itemize}

The \texttt{extra} parameters are passed to the loader function and typically include:
\begin{itemize}
\item \texttt{cols}, \texttt{rows}: Dimensions for reshaping 1D data to 2D
\item \texttt{type}: Data format (\texttt{'2d'}, \texttt{'2drow'})
\item \texttt{delimiter}: For text file parsing
\end{itemize}

\section{Data Loading and Processing}

\subsection{JSON Parsing and Data Loading}

\begin{lstlisting}[language=Python]
def _parse_json(self):
    """Internal method to parse the loaded JSON data using the provided loader function."""
    if not self.json_data:
        raise ValueError("No JSON data loaded. Call load_json() first.")

    # Reset state
    self.frames = {}
    self.extra_data = {}
    self.class_names = []

    # Get the CSV name
    self.feature_table_name = self.json_data.get("csv", "Features_Table.csv")

    # Iterate through the JSON items
    for key, value in self.json_data.items():
        if key == "csv":
            continue

        if isinstance(value, dict):
            class_name = value.get("class", f"Unknown_Class_{key}")
            paths = value.get("path", [])
            extra = value.get("extra", {})

            self.extra_data[class_name] = extra
            self.class_names.append(class_name)
            self.frames[class_name] = []  # Initialize the list for this class

            print(f"[INFO] Loading data for class: {class_name}")
            for path in paths:
                path_obj = Path(path)
                if path_obj.is_file():
                    # USE THE PROVIDED LOADER FUNCTION HERE
                    loaded_data = self.loader_function(path_obj, extra)
                    self._add_loaded_data(loaded_data, class_name)
\end{lstlisting}

\textbf{Explanation:}
This method is the workhorse that processes the JSON configuration:
\begin{enumerate}
\item \textbf{Resets internal state} to ensure clean loading
\item \textbf{Extracts CSV name} for future feature table export
\item \textbf{Iterates through each class definition} in the JSON
\item \textbf{Processes each path} for the class (files or directories)
\item \textbf{Uses the loader function} to actually load the data
\item \textbf{Organizes data by class} in the \texttt{frames} dictionary
\end{enumerate}

The method handles both individual files and directories recursively, making it flexible for various data organization schemes.

\section{Data Balancing Methods}

\subsection{Class Balancing}

\begin{lstlisting}[language=Python]
def balance_frames(self, method='undersample', random_state=None):
    """
    Balance the frames at the data loading level (before feature extraction).
    
    Parameters:
    -----------
    method : str, default='undersample'
        Method to balance classes: 'undersample' or 'oversample'
    
    random_state : int, optional
        Random seed for reproducibility
    """
    if not self.frames:
        raise ValueError("No data loaded. Load data first.")
    
    # Get class distribution
    class_counts = {cls: len(frames) for cls, frames in self.frames.items()}
    print(f"Original class distribution: {class_counts}")
    
    if method == 'undersample':
        self._undersample_frames(random_state)
    elif method == 'oversample':
        self._oversample_frames(random_state)
\end{lstlisting}

\textbf{Explanation:}
Class imbalance is a common problem in machine learning. This method provides two strategies:

\begin{itemize}
\item \textbf{Undersampling}: Reduces samples from majority classes to match the minority class count
\item \textbf{Oversampling}: Increases samples from minority classes using bootstrapping
\end{itemize}

\textbf{Why balance at frame level?}
\begin{itemize}
\item Prevents bias in feature extraction
\item Ensures equal representation during model training
\item More effective than balancing after feature extraction
\end{itemize}

\subsection{Undersampling Implementation}

\begin{lstlisting}[language=Python]
def _undersample_frames(self, random_state=None):
    """Undersample frames to match the minority class"""
    if random_state is not None:
        random.seed(random_state)
        np.random.seed(random_state)
    
    # Find the minority class count
    class_counts = [len(frames) for frames in self.frames.values()]
    minority_count = min(class_counts)
    
    # Undersample each class
    balanced_frames = {}
    for class_name, frames in self.frames.items():
        if len(frames) > minority_count:
            # Undersample
            balanced_frames[class_name] = random.sample(frames, minority_count)
        else:
            # Keep all samples
            balanced_frames[class_name] = frames
    
    self.frames = balanced_frames
\end{lstlisting}

\textbf{Explanation:}
Undersampling works by:
\begin{enumerate}
\item \textbf{Finding the smallest class} (minority count)
\item \textbf{For each class} with more samples than the minority:
  \begin{itemize}
  \item Randomly select a subset of size \texttt{minority\_count}
  \item Uses \texttt{random.sample()} for non-replacement sampling
  \end{itemize}
\item \textbf{Preserves all samples} from classes with fewer samples
\end{enumerate}

\textbf{Advantages:}
\begin{itemize}
\item Faster training (fewer samples)
\item Reduces computational requirements
\item Prevents model from being biased toward majority classes
\end{itemize}

\subsection{Oversampling Implementation}

\begin{lstlisting}[language=Python]
def _oversample_frames(self, random_state=None):
    """Oversample frames to match the majority class"""
    if random_state is not None:
        random.seed(random_state)
        np.random.seed(random_state)
    
    # Find the majority class count
    class_counts = [len(frames) for frames in self.frames.values()]
    majority_count = max(class_counts)
    
    # Oversample each class
    balanced_frames = {}
    for class_name, frames in self.frames.items():
        if len(frames) < majority_count:
            # Oversample with replacement
            balanced_frames[class_name] = resample(
                frames, 
                replace=True, 
                n_samples=majority_count, 
                random_state=random_state
            )
        else:
            # Keep all samples
            balanced_frames[class_name] = frames
    
    self.frames = balanced_frames
\end{lstlisting}

\textbf{Explanation:}
Oversampling works by:
\begin{enumerate}
\item \textbf{Finding the largest class} (majority count)
\item \textbf{For each class} with fewer samples than the majority:
  \begin{itemize}
  \item Uses \texttt{resample()} with replacement to create new samples
  \item Bootstraps existing samples to reach \texttt{majority\_count}
  \end{itemize}
\item \textbf{Preserves all samples} from classes with sufficient samples
\end{enumerate}

\textbf{Advantages:}
\begin{itemize}
\item Preserves all original data
\item Better for small datasets
\item Can improve model performance on minority classes
\end{itemize}

\section{Data Access Methods}

\begin{lstlisting}[language=Python]
def get_data_dict(self):
    """Return the frames dictionary"""
    return self.frames.copy()

def get_class_names(self):
    """Return the list of class names"""
    return self.class_names.copy()

def get_total_samples(self):
    """Return total number of samples across all classes"""
    return sum(len(frames) for frames in self.frames.values())
\end{lstlisting}

\textbf{Explanation:}
These accessor methods provide controlled access to the internal data:
\begin{itemize}
\item \texttt{get\_data\_dict()}: Returns a copy of the frames dictionary to prevent accidental modification
\item \texttt{get\_class\_names()}: Provides the list of unique class labels
\item \texttt{get\_total\_samples()}: Calculates the total dataset size
\end{itemize}

\textbf{Design Philosophy:}
\begin{itemize}
\item \textbf{Encapsulation}: Internal data structures are protected
\item \textbf{Immutability}: Returns copies to prevent external modifications
\item \textbf{Convenience}: Easy access to commonly needed information
\end{itemize}

\chapter{Data Loading Utilities (LovdogData.py)}

The \textbf{LovdogData} module provides specialized loading functions for different data formats, particularly focused on 2D array data commonly used in image processing and computer vision.

\section{Main Loading Function}

\begin{lstlisting}[language=Python]
def load_file_with_extra(file_path, extra_info=None):
    if extra_info is None:
        raise ValueError("extra_info must be provided")

    if 'type' not in extra_info.keys():
        raise ValueError("type must be provided in extra_info")

    if extra_info['type'] == '2d':
        return load_file_with_extra_2d(file_path, extra_info)
    elif extra_info['type'].lower() == '2drow':
        return load_file_with_extra_2dRow(file_path, extra_info)
    else:
        raise ValueError("Valid types are '2d', '2drow', '1drow (not implemented yet)'")
\end{lstlisting}

\textbf{Explanation:}
This function serves as a dispatcher that routes to the appropriate loading function based on the data type specified in \texttt{extra\_info}:

\begin{itemize}
\item \texttt{'2d'}: For files containing a single 2D array
\item \texttt{'2drow'}: For files where each line represents a flattened 2D array
\item \texttt{'1drow'}: Reserved for future 1D data implementations
\end{itemize}

\textbf{Key Features:}
\begin{itemize}
\item \textbf{Type safety}: Validates that \texttt{extra\_info} contains required parameters
\item \textbf{Extensibility}: Easy to add new data format handlers
\item \textbf{Error handling}: Clear error messages for missing or invalid parameters
\end{itemize}

\section{2D Array Loading}

\begin{lstlisting}[language=Python]
def load_file_with_extra_2d(file_path, extra_info=None):
    """
    Loads a file containing a single 2D array.
    """
    if extra_info is None:
        extra_info = {}
    
    file_path = Path(file_path)
    if not file_path.exists():
        print(f"Error: File not found: {file_path}")
        return []

    # Get the essential shape information
    try:
        cols = int(extra_info['cols'])
        rows = int(extra_info['rows'])
    except KeyError:
        print(f"Error: 'cols' and 'rows' must be provided in extra_info for file {file_path}.")
        return []
    
    frames_per_file = []
    delimiter = extra_info.get('delimiter', None)
    
    try:
        frame_2d = np.loadtxt(file_path, delimiter=delimiter).reshape((rows, cols))
        frames_per_file.append(frame_2d)
                
        print(f"DEBUG: Loaded {len(frames_per_file)} frames from {file_path}")
        return frames_per_file
\end{lstlisting}

\textbf{Explanation:}
This function loads files that contain a single 2D array:

\begin{enumerate}
\item \textbf{File validation}: Checks if the file exists
\item \textbf{Parameter extraction}: Gets dimensions from \texttt{extra\_info}
\item \textbf{Data loading}: Uses \texttt{np.loadtxt()} to read the file
\item \textbf{Reshaping}: Converts flattened data to 2D using specified dimensions
\item \textbf{Return format}: Always returns a list (for consistency)
\end{enumerate}

\textbf{Typical Use Cases:}
\begin{itemize}
\item Single images stored as text files
\item Pre-computed feature matrices
\item 2D sensor data readings
\end{itemize}

\section{2D Row-wise Loading}

\begin{lstlisting}[language=Python]
def load_file_with_extra_2dRow(file_path, extra_info=None):
    """
    Loads a file where each line is a flattened 1D array representing a single frame.
    """
    # [Parameter validation similar to above]
    
    try:
        with open(file_path, 'r') as f:
            for line_num, line in enumerate(f):
                # Skip empty lines
                if not line.strip():
                    continue
                
                # Convert the line of text into a list of numbers
                data_1d = np.loadtxt([line], delimiter=delimiter)
                
                # Check if the line has the correct number of elements
                expected_elements = cols * rows
                if data_1d.size != expected_elements:
                    print(f"Warning: Line {line_num+1} has {data_1d.size} elements, "
                          f"but {expected_elements} are needed for a {rows}x{cols} frame. Skipping line.")
                    continue
                
                # Reshape the 1D array into a 2D frame
                frame_2d = data_1d.reshape((rows, cols))
                frames_per_file.append(frame_2d)
\end{lstlisting}

\textbf{Explanation:}
This function handles files where each line represents a separate 2D array:

\begin{enumerate}
\item \textbf{Line-by-line processing}: Reads each line independently
\item \textbf{Data validation}: Checks if each line has the correct number of elements
\item \textbf{Individual reshaping}: Converts each line to a 2D array
\item \textbf{Robust error handling}: Skips invalid lines but continues processing
\end{enumerate}

\textbf{Typical Use Cases:}
\begin{itemize}
\item Time series of 2D data (video frames, sensor sequences)
\item Multiple observations in a single file
\item Batch processing of related 2D arrays
\end{itemize}

\textbf{Key Differences from \texttt{load\_file\_with\_extra\_2d}:}
\begin{itemize}
\item Returns \textbf{multiple frames} (one per valid line)
\item Handles \textbf{mixed valid/invalid} data gracefully
\item Suitable for \textbf{sequential data} processing
\end{itemize}

\section{Data Augmentation}

\begin{lstlisting}[language=Python]
def translate_px_aug(px, extra_info):
    if extra_info is None:
        extra_info = {}

    offset_x = extra_info.get('offset_x', 1)
    offset_y = extra_info.get('offset_y', 1)
    tremble  = extra_info.get('tremble', True)
    
    px_new = px.copy()  # Start with a copy of the original list

    for p in px:
        # Translate the pixels positively
        px_mod = np.roll(p, shift=offset_y, axis=0)
        px_mod = np.roll(px_mod, shift=offset_x, axis=1)
        np.append(px_new, px_mod)

    if tremble:
        for p in px:
            # Translate the pixels negatively
            px_mod = np.roll(p, shift=-offset_y, axis=0)
            px_mod = np.roll(px_mod, shift=-offset_x, axis=1)
            np.append(px_new, px_mod)

    return np.array(px_new)
\end{lstlisting}

\textbf{Explanation:}
This function provides data augmentation through pixel translation:

\begin{itemize}
\item \textbf{Positive translation}: Shifts pixels right/down using \texttt{np.roll()}
\item \textbf{Negative translation}: Shifts pixels left/up (when \texttt{tremble=True})
\item \textbf{Circular boundary handling}: \texttt{np.roll()} wraps around edges
\item \textbf{Parameter control}: Translation amounts configurable via \texttt{extra\_info}
\end{itemize}

\textbf{Augmentation Benefits:}
\begin{itemize}
\item \textbf{Increased dataset size} without collecting new data
\item \textbf{Improved model robustness} to spatial variations
\item \textbf{Reduced overfitting} by presenting varied examples
\end{itemize}

\textbf{Use with LovdogDataFrames:}
\begin{lstlisting}[language=Python]
# Apply augmentation to loaded data
data_loader.augmentation_data(translate_px_aug, 
                             {'offset_x': 2, 'offset_y': 2, 'tremble': True})
\end{lstlisting}

\chapter{Feature Extraction (LovdogLBP.py)}

The \textbf{LBPFeatures} class implements Local Binary Pattern feature extraction, a powerful texture analysis technique commonly used in image processing and computer vision applications.

\section{Main Class: LBPFeatures}

\begin{lstlisting}[language=Python]
class LBPFeatures:
    features = [
        "File",
        "mean",
        "variance", 
        "correlation",
        "contrast",
        "homogeneity",
        "class"
    ]
    verbose = 1

    def __init__(self, data_loader=None):
        self.data_loader = data_loader
        self.LBP = {}
        self.features_table = pd.DataFrame(columns=self.features)
        
        # LBP parameters with defaults
        self.ratio = 1
        self.points = 8
        self.method = "uniform"
        self.window_size = [3]
\end{lstlisting}

\textbf{Explanation:}
The LBPFeatures class extracts texture features using the Local Binary Pattern algorithm:

\begin{itemize}
\item \texttt{features}: Defines the output feature columns (statistical measures)
\item \texttt{data\_loader}: Reference to a LovdogDataFrames instance for data access
\item \texttt{LBP}: Stores intermediate LBP images for each class
\item \texttt{features\_table}: Final DataFrame containing extracted features
\end{itemize}

\textbf{LBP Parameters:}
\begin{itemize}
\item \texttt{points}: Number of circularly symmetric neighbor points (typically 8, 16, or 24)
\item \texttt{ratio}: Radius of the circle (scaled by points)
\item \texttt{method}: LBP variant ('default', 'ror', 'uniform', 'var')
\item \texttt{window\_size}: For sliding window analysis of LBP images
\end{itemize}

\section{LBP Computation}

\begin{lstlisting}[language=Python]
def compute_lbp_from_data_loader(self):
    """
    Compute LBP features using data from the LovdogDataFrames object.
    """
    if self.data_loader is None:
        raise ValueError("No data loader provided. Call set_data_loader() first.")
    
    frames_dict = self.data_loader.get_data_dict()
    class_names = self.data_loader.get_class_names()

    for class_name in class_names:
        print(f"  Processing class: {class_name}")
        self.LBP[class_name] = []
        
        for i, frame in enumerate(frames_dict[class_name]):
            # Convert frame to uint8 if needed
            if isinstance(frame, np.ndarray):
                # Normalize and convert to uint8 for LBP processing
                frame_normalized = cv.normalize(frame, None, 0, 255, cv.NORM_MINMAX).astype(np.uint8)
                lbp = local_binary_pattern(frame_normalized, self.points, self.ratio, self.method).astype(np.uint8)
                
                # Use a generic filename or index
                file_id = f"{class_name}_{i:04d}"
                self.LBP[class_name].append({"File": file_id, "LBP": lbp, "class": class_name})
\end{lstlisting}

\textbf{Explanation:}
This method processes data from the data loader through the LBP pipeline:

\begin{enumerate}
\item \textbf{Data access}: Retrieves frames organized by class from the data loader
\item \textbf{Normalization}: Converts input data to 8-bit format required by LBP
\item \textbf{LBP computation}: Applies the Local Binary Pattern algorithm
\item \textbf{Result storage}: Stores LBP images with metadata for later feature extraction
\end{enumerate}

\textbf{LBP Algorithm Overview:}
\begin{itemize}
\item \textbf{Local analysis}: Examines each pixel relative to its neighbors
\item \textbf{Binary patterns}: Creates binary codes representing local texture
\item \textbf{Rotation invariance}: 'uniform' method provides rotation-invariant features
\item \textbf{Texture encoding}: Converts local texture patterns into numerical codes
\end{itemize}

\section{Feature Extraction from LBP}

\begin{lstlisting}[language=Python]
def compute_features_from_lbp(self):
    """Compute features from LBP images (works with both old and new data sources)"""
    bin_total = 25
    bin_value = 0
    for c in self.LBP.keys():
        print(f"Processing {c} images...")
        fd_size = len(self.LBP[c])
        bin_value = fd_size // bin_total
        fd_count = 0
        for i in self.LBP[c]:
            for ws in self.window_size:
                result = self.lbp_window_analysis(i, ws, _class=c)
                self.features_table = LBPFeatures.append_features_row(result, self.features_table)
            if fd_count % bin_value == 0:
                print("*", end="", flush=True)
            fd_count += 1
\end{lstlisting}

\textbf{Explanation:}
This method converts LBP images into statistical features:

\begin{itemize}
\item \textbf{Progress tracking}: Visual feedback during processing
\item \textbf{Window analysis}: Applies sliding window to extract local features
\item \textbf{Multiple scales}: Processes each window size specified
\item \textbf{Feature aggregation}: Combines results into final feature table
\end{itemize}

\textbf{Progress Indicators:}
\begin{itemize}
\item \texttt{bin\_total}: Number of progress markers to display
\item \texttt{bin\_value}: Calculates how many processed per marker
\item Visual feedback helps monitor long-running feature extraction
\end{itemize}

\section{Window-based Feature Analysis}

\begin{lstlisting}[language=Python]
@staticmethod
def lbp_window_analysis(lbp_img, window_size, _class=None):
    x = 0
    y = 0
    lbp_image = lbp_img["LBP"]
    lbp_file = lbp_img["File"]
    analysis = []
    run_boggle = [ "-", "\\", "|", "/" ]
    idx_boggle = 0

    if window_size < 0:
        # Process entire image as single window
        lbp_feat = LBPFeatures.get_features_from_window(lbp_image)
        lbp_feat["File"] = lbp_file
        if _class is not None:
            lbp_feat["class"] = _class
        analysis.append(lbp_feat)
        return pd.DataFrame(analysis, columns=LBPFeatures.features)
    
    # Sliding window approach
    for x in range(lbp_image.shape[0] - window_size):
        for y in range(lbp_image.shape[1] - window_size):
            window = lbp_image[x:x + window_size, y:y + window_size]
            lbp_feat = LBPFeatures.get_features_from_window(window)
            lbp_feat["File"] = lbp_file
            if _class is not None:
                lbp_feat["class"] = _class
            analysis.append(lbp_feat)
            # Progress animation
            print(run_boggle[idx_boggle], end="", flush=True)
            idx_boggle = idx_boggle + 1 if idx_boggle < 3 else 0
            print("\b", end="", flush=True)
\end{lstlisting}

\textbf{Explanation:}
This method implements sliding window analysis of LBP images:

\begin{itemize}
\item \textbf{Two modes}: Whole-image analysis (window\_size < 0) or sliding window
\item \textbf{Sliding window}: Extracts features from overlapping sub-regions
\item \textbf{Progress animation}: Rotating cursor shows processing activity
\item \textbf{Feature aggregation}: Collects features from all window positions
\end{itemize}

\textbf{Window Analysis Benefits:}
\begin{itemize}
\item \textbf{Spatial information}: Captures local texture variations
\item \textbf{Multi-scale analysis}: Different window sizes capture different texture scales
\item \textbf{Robust features}: Averages out local noise through multiple observations
\end{itemize}

\section{GLCM Feature Extraction}

\begin{lstlisting}[language=Python]
@staticmethod
def get_features_from_window(window):
    features = {}
    glcm = graycomatrix(window, [1], [0], symmetric=True, normed=True, levels=59)
    features["mean"] = graycoprops(glcm, 'mean')[0][0]
    features["variance"] = graycoprops(glcm, 'variance')[0][0]
    features["correlation"] = graycoprops(glcm, 'correlation')[0][0]
    features["contrast"] = graycoprops(glcm, 'contrast')[0][0]
    features["homogeneity"] = graycoprops(glcm, 'homogeneity')[0][0]
    return features
\end{lstlisting}

\textbf{Explanation:}
This static method computes Gray-Level Co-occurrence Matrix (GLCM) features from LBP windows:

\begin{itemize}
\item \textbf{GLCM computation}: Creates a matrix showing how pixel pairs co-occur
\item \textbf{Parameter choices}:
  \begin{itemize}
  \item \texttt{distance=1}: Neighboring pixels
  \item \texttt{angle=0}: Horizontal direction
  \item \texttt{symmetric=True}: Reduces computation and provides rotational invariance
  \item \texttt{levels=59}: Quantization levels for LBP values (8-bit range)
  \end{itemize}
\item \textbf{Statistical features}: Five key texture measures
\end{itemize}

\textbf{GLCM Feature Definitions:}
\begin{itemize}
\item \textbf{Mean}: Average gray level intensity
\item \textbf{Variance}: Spread of gray level distribution
\item \textbf{Correlation}: Linear dependency of gray levels
\item \textbf{Contrast}: Local intensity variations (edge strength)
\item \textbf{Homogeneity}: Closeness of element distribution to GLCM diagonal
\end{itemize}

\section{Machine Learning Data Preparation}

\begin{lstlisting}[language=Python]
def get_ml_data(self, include_file=False):
    """
    Return features and labels in standard scikit-learn format.
    
    Args:
        include_file (bool): Whether to include the 'File' column in features
    
    Returns:
        tuple: (X_features, y_labels) where:
               X_features: numpy array or DataFrame of features
               y_labels: numpy array or Series of class labels
    """
    if include_file:
        X = self.features_table.drop('class', axis=1)
    else:
        X = self.features_table.drop(['File', 'class'], axis=1)
    
    y = self.features_table['class']
    return X, y
\end{lstlisting}

\textbf{Explanation:}
This method converts the feature table into standard machine learning format:

\begin{itemize}
\item \texttt{include\_file=False}: Standard case, returns only numerical features
\item \texttt{include\_file=True}: Includes filename for debugging or tracking
\item \textbf{X}: Feature matrix (samples × features)
\item \textbf{y}: Label vector (samples × 1)
\end{itemize}

\textbf{Compatibility:}
\begin{itemize}
\item Direct input to scikit-learn classifiers
\item Compatible with train/test splitting
\item Works with cross-validation and grid search
\end{itemize}

\chapter{Model Selection (LovdogMS.py)}

The \textbf{ModelSelection} class is the core of the automated machine learning framework, providing comprehensive model evaluation, hyperparameter tuning, and performance analysis.

\section{Main Class: ModelSelection}

\begin{lstlisting}[language=Python]
class ModelSelection:
    def __init__(
        self, 
        name: str, 
        # Multiple data input options for flexibility
        csv_path: Optional[str] = None,
        features_dataframe: Optional[pd.DataFrame] = None,
        ml_data_tuple: Optional[Tuple] = None,
        # Dataset configuration
        target_column: str = "class",
        test_size: float = 0.2,
        random_state: Optional[int] = None,
        # Performance optimization
        best_only_mode: bool = False,
        # Output configuration
        plots_output_directory: str = "plots",
        # Classifier configuration
        classifiers_dict: Optional[Dict[str, Any]] = None,
        grid_parameters_dict: Optional[Dict[str, Any]] = None
    ):
\end{lstlisting}

\textbf{Explanation:}
The ModelSelection class provides a unified interface for automated model evaluation:

\begin{itemize}
\item \texttt{name}: Unique identifier for tracking multiple experiments
\item \textbf{Multiple data inputs}: CSV files, DataFrames, or direct (X,y) tuples
\item \texttt{target\_column}: Specifies which column contains class labels
\item \texttt{test\_size}: Train/test split ratio (0.2 = 80\% train, 20\% test)
\item \texttt{best\_only\_mode}: Optimizes performance by only analyzing the best model
\item \texttt{plots\_output\_directory}: Where to save diagnostic plots
\end{itemize}

\textbf{Data Input Flexibility:}
\begin{itemize}
\item \texttt{csv\_path}: Direct file loading for quick experiments
\item \texttt{features\_dataframe}: Pre-loaded pandas DataFrame
\item \texttt{ml\_data\_tuple}: Direct (X,y) input from LBP feature extraction
\end{itemize}

\section{Default Classifier Configuration}

\begin{lstlisting}[language=Python]
default_classifiers = {
    'SVC': SVC(),
    'KNeighbors': KNeighborsClassifier(),
    'RandomForest': RandomForestClassifier(),
    'LogisticRegression': LogisticRegression(max_iter=850),
}

default_grid_parameters = {
    'SVC': {
        'clf__C': [0.001, 0.01, 0.1, 1, 10, 20],
        'clf__kernel': ['poly', 'sigmoid', 'rbf'],
        'clf__gamma': ['scale', 'auto'],
    },
    'RandomForest': {
        'clf__min_samples_split': [2, 4, 8, 16],
        'clf__n_estimators': [72, 144, 288, 576],
        'clf__max_depth': [3, 9, 27],
    }
}
\end{lstlisting}

\textbf{Explanation:}
The framework includes four well-established classifiers with optimized parameter grids:

\begin{itemize}
\item \textbf{Support Vector Classifier (SVC)}: Effective for high-dimensional data
\item \textbf{K-Nearest Neighbors}: Simple, interpretable, good for small datasets
\item \textbf{Random Forest}: Robust, handles non-linear relationships well
\item \textbf{Logistic Regression}: Fast, good baseline, probabilistic outputs
\end{itemize}

\textbf{Parameter Grid Design:}
\begin{itemize}
\item \textbf{Exponential scales}: For regularization parameters (0.001 to 20)
\item \textbf{Power-of-two progressions}: For ensemble sizes (72 to 576)
\item \textbf{Multiple strategies}: Different kernel types, splitting criteria
\item \textbf{Computational efficiency}: Balanced between exploration and computation time
\end{itemize}

\section{Core Model Selection Process}

\begin{lstlisting}[language=Python]
def selectModels(self):
    """Run grid search for all classifiers, including AUC calculation"""
    self.results = {}
    all_metrics = []
    
    for tag, algorithm in self.classifiers.items():
        print(f"Training {tag}...")
        
        # Prepare pipeline
        pipe = Pipeline([
            ('scaler', StandardScaler()),
            ('clf', algorithm)
        ])
        
        # SetUp grid
        modelsGrid = GridSearchCV(
            pipe,
            param_grid=self.gridParameters[tag],
            scoring={
                "accuracy": "accuracy",
                "precision": "precision_weighted",
                "recall": "recall_weighted", 
                "f1": "f1_weighted"
            },
            refit="accuracy",
            cv=5,
            n_jobs=-1,
            verbose=1
        )
\end{lstlisting}

\textbf{Explanation:}
The main model selection method implements a comprehensive evaluation pipeline:

\begin{enumerate}
\item \textbf{Pipeline construction}: StandardScaler + Classifier for consistent preprocessing
\item \textbf{GridSearchCV setup}: Exhaustive hyperparameter search with cross-validation
\item \textbf{Multi-metric evaluation}: Accuracy, Precision, Recall, F1-score
\item \textbf{Parallel processing}: \texttt{n\_jobs=-1} uses all available CPU cores
\item \textbf{5-fold cross-validation}: Robust performance estimation
\end{enumerate}

\textbf{Pipeline Benefits:}
\begin{itemize}
\item \textbf{Automatic scaling}: Features standardized before classification
\item \textbf{Prevention of data leakage}: Scaling fitted only on training folds
\item \textbf{Consistent preprocessing}: Same transformation for all models
\end{itemize}

\section{Comprehensive Performance Metrics}

\begin{lstlisting}[language=Python]
# Calculate all basic metrics
accuracy = accuracy_score(self.test_data['y'], y_pred)
precision = precision_score(self.test_data['y'], y_pred, average='weighted')
recall = recall_score(self.test_data['y'], y_pred, average='weighted')
f1 = f1_score(self.test_data['y'], y_pred, average='weighted')

# Calculate AUC if the model supports probability predictions
auc_score = None
if hasattr(best_model, 'predict_proba'):
    try:
        y_pred_proba = best_model.predict_proba(self.test_data['x'])
        
        # For multi-class classification, use One-vs-Rest approach
        if len(np.unique(self.test_data['y'])) > 2:
            y_test_bin = label_binarize(self.test_data['y'], 
                                      classes=np.unique(self.test_data['y']))
            # Calculate micro-average AUC
            fpr, tpr, _ = roc_curve(y_test_bin.ravel(), y_pred_proba.ravel())
            auc_score = auc(fpr, tpr)
        else:
            # For binary classification
            fpr, tpr, _ = roc_curve(self.test_data['y'], y_pred_proba[:, 1])
            auc_score = auc(fpr, tpr)
\end{lstlisting}

\textbf{Explanation:}
The framework computes multiple performance metrics for comprehensive model evaluation:

\begin{itemize}
\item \textbf{Accuracy}: Overall correct classification rate
\item \textbf{Precision}: Ratio of true positives to predicted positives
\item \textbf{Recall}: Ratio of true positives to actual positives  
\item \textbf{F1-score}: Harmonic mean of precision and recall
\item \textbf{AUC-ROC}: Area Under ROC Curve for probabilistic models
\end{itemize}

\textbf{Metric Selection Rationale:}
\begin{itemize}
\item \textbf{Weighted averages}: Handle multi-class and imbalanced datasets
\item \textbf{AUC calculation}: Only for models supporting probability predictions
\item \textbf{Micro-averaging}: For multi-class AUC computation
\item \textbf{Comprehensive view}: Different metrics highlight different strengths
\end{itemize}

\section{Best Model Selection}

\begin{lstlisting}[language=Python]
# Determine the best classifier (considering AUC if available)
if not self.metrics_df.empty:
    # Prefer classifiers that have AUC scores for ranking
    has_auc_df = self.metrics_df[self.metrics_df['AUC_Score'] != 'N/A'].copy()
    if not has_auc_df.empty:
        # Convert AUC_Score to float for sorting
        has_auc_df['AUC_Score'] = has_auc_df['AUC_Score'].astype(float)
        best_idx = has_auc_df['AUC_Score'].idxmax()
        self.best_model_tag = has_auc_df.loc[best_idx, 'Classifier']
    else:
        # Fallback to accuracy if no AUC available
        best_idx = self.metrics_df['Accuracy'].idxmax()
        self.best_model_tag = self.metrics_df.loc[best_idx, 'Classifier']
    
    self.best_model = self.results[self.best_model_tag]['best_estimator']
\end{lstlisting}

\textbf{Explanation:}
The framework implements a smart best-model selection strategy:

\begin{enumerate}
\item \textbf{Priority to AUC}: If available, AUC is preferred for model ranking
\item \textbf{Fallback to accuracy}: For models without probability outputs
\item \textbf{Automatic type conversion}: Handles string 'N/A' values gracefully
\item \textbf{Model retrieval}: Stores both the best model and its metadata
\end{enumerate}

\textbf{Why Prefer AUC?}
\begin{itemize}
\item \textbf{Threshold-independent}: Not affected by classification threshold choice
\item \textbf{Probability quality}: Measures how well probabilities are calibrated
\item \textbf{Robust to imbalance}: Works well with unbalanced class distributions
\end{itemize}

\section{Diagnostic Visualization}

\begin{lstlisting}[language=Python]
def generatePlots(self):
    """Generate plots - only for best classifier if best_only=True, otherwise for all"""
    if self.best_only_mode:
        self._generatePlotsForBest()
    else:
        self._generatePlotsForAll()

def _generateLearningCurve(self, tag):
    """Generate learning curve for a specific classifier"""
    model = self.results[tag]['best_estimator']
    
    train_sizes, train_scores, test_scores = learning_curve(
        model, self.data, self.target, cv=5, n_jobs=-1,
        train_sizes=np.linspace(0.1, 1.0, 5),
        scoring='accuracy'
    )
\end{lstlisting}

\textbf{Explanation:}
The framework generates three types of diagnostic plots:

\begin{itemize}
\item \textbf{Learning Curves}: Show how performance changes with training set size
\item \textbf{ROC Curves}: Display trade-off between true positive and false positive rates
\item \textbf{Confusion Matrices}: Detailed breakdown of classification errors
\end{itemize}

\textbf{Plot Generation Modes:}
\begin{itemize}
\item \texttt{best\_only\_mode=True}: Only generate plots for the best-performing model
\item \texttt{best\_only\_mode=False}: Generate plots for all evaluated models
\end{itemize}

\textbf{Learning Curve Interpretation:}
\begin{itemize}
\item \textbf{Large gap}: Indicates overfitting (model too complex)
\item \textbf{Convergence}: Both curves plateau at similar levels (good fit)
\item \textbf{Poor performance}: Both curves remain low (underfitting)
\end{itemize}

\section{Model Persistence and Export}

\begin{lstlisting}[language=Python]
def saveModels(self, path):
    """Save all trained models to disk"""
    os.makedirs(path, exist_ok=True)
    for tag, result in self.results.items():
        joblib.dump(result['best_estimator'], f"{path}/{tag}.joblib")

def getMetricsDataFrame(self):
    """Create a DataFrame with all metrics including AUC for CSV export"""
    if self.metrics_df.empty:
        # If selectModels hasn't been run, create basic metrics dataframe
        metrics_data = []
        for tag, result in self.results.items():
            metrics_data.append({
                'Classifier': tag,
                'Accuracy': result['test_accuracy'],
                'Precision': result['test_precision'],
                'Recall': result['test_recall'],
                'F1_Score': result['test_f1'],
                'AUC_Score': result['test_auc'] if result['test_auc'] is not None else 'N/A',
                'CV_Score': result['cv_best_score'],
                'Best_Params': str(result['best_params']),
                'Supports_Probability': result['supports_proba']
            })
        
        self.metrics_df = pd.DataFrame(metrics_data)
    
    return self.metrics_df
\end{lstlisting}

\textbf{Explanation:}
The framework provides comprehensive model and results export:

\begin{itemize}
\item \textbf{Model serialization}: All trained models saved using joblib
\item \textbf{Directory creation}: Automatic creation of output directories
\item \textbf{Comprehensive metrics}: Detailed performance statistics in DataFrame format
\item \textbf{CSV export ready}: Direct compatibility with spreadsheet software
\item \textbf{Parameter documentation}: Best hyperparameters included for reproducibility
\end{itemize}

\textbf{Export Benefits:}
\begin{itemize}
\item \textbf{Reproducibility}: All models and parameters saved
\item \textbf{Analysis ready}: Metrics formatted for further statistical analysis
\item \textbf{Deployment ready}: Trained models can be loaded for inference
\item \textbf{Documentation}: Complete record of experimental results
\end{itemize}

\end{document}
